\appendixheading

*引理 1. --- 令 X 为一个局部紧空间，R 为 X 上的开等价关系，且商空间 X/R 为仿紧；令 \(\pi\) 为从 X 到 X/R 的典范映射。存在 X 上一个连续的实值函数 \(F \geqslant 0\)，使得：

\begin{enumerate}
\def\labelenumi{\alph{enumi})}
\item
  F 关于 R 的每个等价类都不恒为零；
\item
  对于 X/R 的每个紧子集 K，\(\overline{\pi}^{1}(K)\) 与 Supp F 的交是紧的。
\end{enumerate}

对每个点 \(u \in \mathbf{X} / \mathbf{R}\)，选取一个函数 \(f_{u} \in \mathcal{K}_{+}(\mathbf{X})\)，使得 \(f_{u}\) 在 \(\overline{\pi}(u)\) 上不恒为零；令 \(\Omega_{u}\) 为 \(f_{u} > 0\) 的点所成的开集；于是 \(u \in \pi(\Omega_{u})\)。由于 \(\pi\) 是开映射，\(\pi(\Omega_{u})\) 构成 \(\mathbf{X} / \mathbf{R}\) 的一个开覆盖。存在一个局部有限开覆盖 \((\mathrm{U}_{\iota})_{\iota \in \mathbf{I}}\)，它比由 \(\pi(\Omega_{u})\) 构成的覆盖更细，于是（GT, IX, §4, No.~3, Prop. 3）存在单位分解 \((g_{\iota})_{\iota \in \mathbf{I}}\)，它定义在 \(\mathbf{X} / \mathbf{R}\) 上并从属于覆盖 \((\mathrm{U}_{\iota})\)。对每个 \(\iota \in \mathbf{I}\)，选取 \(u_{\iota}\) 使得 \(\mathrm{U}_{\iota} \subset \pi(\Omega_{u_{\iota}})\)。函数 \(\mathrm{F}_{\iota} = (g_{\iota} \circ \pi) \cdot f_{u_{\iota}}\) 属于 \(\mathcal{K}(\mathbf{X})\)，且其支集包含于 \(\overline{\pi}(\mathrm{U}_{\iota})\) 中。因此，\(F_{\iota}\) 的支集构成局部有限族，从而可以定义连续函数 \(F \geqslant 0\)，它在 \(X\) 上：令 \(F = \sum_{\iota \in I} F_{\iota}\)。

对于每个 \(u \in X/R\)，存在某个 \(\iota\) 使得 \(g_{\iota}(u) > 0\)，因而 \(u \in U_{\iota}\)；接着存在 \(x \in \Omega_{u_{\iota}}\) 使得 \(\pi(x) = u\)；于是 \(f_{u_{\iota}}(x) > 0\) 且 \(g_{\iota}(\pi(x)) > 0\)，所以 \(F_{\iota}(x) > 0\)，从而 \(F(x) > 0\)；这证明了 F 具有性质 a)。最后，令 K 为 X/R 的一个紧子集。存在 I 的有限子集 J，使得对于 \(\iota \in I - J\)，有 \(\mathrm{U}_{\iota}\cap \mathrm{K} = \varnothing\)，因而 \(\overline{\pi} (\mathbf{K})\cap \operatorname {Supp}\mathbf{F}_{\iota} = \varnothing\)。于是集合

\[
\overline {{\pi}} ^ {1} (\mathrm{K}) \cap \operatorname{Supp} \mathrm{F} = \overline {{\pi}} ^ {1} (\mathrm{K}) \cap \left(\bigcup_ {\iota \in \mathrm{I}} \operatorname{Supp} \mathrm{F} _ {\iota}\right) = \overline {{\pi}} ^ {1} (\mathrm{K}) \cap \left(\bigcup_ {\iota \in \mathrm{J}} \operatorname{Supp} \mathrm{F} _ {\iota}\right)
\]

是紧的。

引理 2. --- 设 G 为在无穷远处可数的局部紧群，M 为 Baire 空间。假设 G 在 M 上连续且传递地左作用。对每个 \(x \in M\)，令 \(H_{x}\) 为 x 在 G 中的稳定子群，于是映射 \(s \mapsto sx\) 诱导出一个连续双射 \(\varphi_{x}\)，将 \(G/H_{x}\) 映到 M。那么 \(\varphi_{x}\) 是从 \(G/H_{x}\) 到 M 的同胚（换言之（GT, III, §2, No.~5），M 是一个拓扑齐性空间）。

令 \(x_0 \in \mathbf{M}\)。只需证明（同上，Prop. 15），映射 \(s \mapsto sx_0\) 将每个邻域 \(\mathbf{V}\)（其中 \(e\) 在 \(\mathbf{G}\) 中）变为 \(x_0\) 在 \(\mathbf{M}\) 中的一个邻域。令 \(\mathbf{W}\) 为 \(e\) 的对称紧邻域，使得 \(\mathbf{W}^2 \subset \mathbf{V}\)。根据假设，\(\mathbf{G}\) 是一列紧集的并，因而是平移列 \((s_n\mathbf{W})\)（它们是 \(\mathbf{W}\) 的平移）的并。于是 \(\mathbf{M}\) 是一列紧集 \((s_n\mathbf{W}x_0)\) 的并。由于 \(\mathbf{M}\) 是 Baire 空间，存在某个 \(n\)，使 \(s_n\mathbf{W}x_0\) 有一个内点 \(s_nwx_0\)（其中 \(w \in \mathbf{W}\)）。因此 \(x_0\) 是

\[
w ^ {- 1} s _ {n} ^ {- 1} (s _ {n} \mathrm{W} x _ {0}) = w ^ {- 1} \mathrm{W} x _ {0} \subset \mathrm{V} x _ {0},
\]

从而 \(Vx_{0}\) 是 M 中 \(x_{0}\) 的一个邻域。

